\chapter{训练数据工程}
\label{lab:20_dataset}

\noindent\textbf{配套文件：}\path{notebook/20_dataset.ipynb}\par
\noindent\textbf{教材关联：}\bookxrefshort{ch:data-engineering}、\bookxrefshort{ch:pretraining-scaling}、\bookxrefshort{ch:supervised-finetuning}。

\section*{实验目标}
把记录、Schema、内容身份和划分规则共同视为数据资产。

\section*{准备及定位}
先阅读本册实验总则，确认Notebook中的依赖与资产单元。原理计算、库接口迁移和真实模型训练可能具有不同资源要求，应分别记录设备、版本及运行范围。

源码定位可从下列函数开始；应结合前置数据与配置单元阅读。
\begin{itemize}
\item \texttt{my\_locate\_project\_directory}
\item \texttt{my\_build\_parallel\_pairs}
\item \texttt{my\_add\_quality\_columns}
\end{itemize}

\section*{操作及观察}
\begin{enumerate}
\item 逐条检查候选句对的字段、空值与唯一身份，保留被拒记录的原因。
\item 执行JSONL序列化与原子发布，核对重载记录数及内容哈希。
\item 通过Datasets加载同一权威源，比较map/filter前后的记录保留关系。
\item 生成稳定划分并持久化Arrow，重载后验证Schema、样本身份和划分交集。
\end{enumerate}

\section*{结果判读及提交}
提交数据清单、哈希、划分交集和重载对照。文件大小相同不证明内容相同；新增语料后须重新审查泄漏，不能只验证行数。

提交时附上所执行单元范围、环境与制品身份、原始结果及失败说明。将未运行项目明确列出，不能用Notebook中已有输出替代本次执行证据。
